\documentclass{article}

% Language setting
% Replace `english' with e.g. `spanish' to change the document language
\usepackage[english]{babel}

% Set page size and margins
% Replace `letterpaper' with `a4paper' for UK/EU standard size
\usepackage[letterpaper,top=2cm,bottom=2cm,left=3cm,right=3cm,marginparwidth=1.75cm]{geometry}

% Useful packages
\usepackage{amsmath}
\usepackage{graphicx}
\usepackage[colorlinks=true, allcolors=blue]{hyperref}

\title{Interests in Data Science}
\author{Caleb Reid}

\begin{document}
\maketitle
\section{Prompt}
 What are you  your interests in economics and data science. What made you want to take this class?
Do you have any ideas for what you would want to do for your project for this class? 
What are your goals for this class, and what is your plan for after graduation? 

\section{Response}
I first got interested in data science and analytics during my second semester of sophomore year. I took a class called sports analytics and realized I had a passion for analytics. While science was never my favorite subject, I've always been interested in all of the stats surrounding sports. In fact, my mom was the first person to tell me I should be doing something related to sports analytics. I never thought it would become something I pursue a career in, but I have found I really enjoy the science behind data analytics.

I took this class because my professor highly recommended both the course and the professor. It aligns with my degree, sports data analytics, and I expect to learn a lot. I want to become better at organizing my files and writing clean code. I hope that I am able to choose a topic for my thesis proposal and get some initial work done before the end of the semester. I plan to pursue a career in sports analytics post-graduation. I would prefer to work in professional field for sports like football and soccer, but I am open to whatever opportunities are out there. I also have an undergrad certificate in esports business and analytics, which gives me additional options post-graduation. 

\section{Equation}

\(a^2 + b^2 = c^2\)

\end{document}
